\section{Prym 写像の像と逆構成を使うためのデータ}
\label{sec:notes-prym-v27}

主要な分類の結論は次のように分かれる。ここで $\mathcal R_g$ は
曲線と非自明二捩れ線束の対のモジュライであり、Prym 多様体のモジュライではない。
\begin{center}\small
\begin{tabular}{p{.15\linewidth}p{.20\linewidth}p{.53\linewidth}}\toprule
署名と次元&像&復元と残るファイバー\\\midrule
$(2,4)$、八次元&$\mathcal R_3$ を支配&一般関数体上で $E\times D$ と双有理。$E$ の周期・指数は三、$D$ は非分裂二次曲線、ファイバーの指数・最小多重切断次数は六。\\
$(3,3)$、七次元&$\mathcal R_{10}$ 内の七次元像&抽象 $(C,\eta)$ から双有理復元。$W_6^2(C)=\{H\}$ は被約。指定反対合を保てば固有ブロックの順序も不要。\\\bottomrule
\end{tabular}\end{center}
幾何と六次の復元は公表稿V、net写像の相対有理性と一般関数体上の算術は公表稿VIに全証明を置く。
指定反対合そのものを忘れた代数の粗モジュライへの写像は、この表の結論に含めない。

\subsection{全incidenceから係数Gram行列を読む}
$P$ の行数を $p$ とし、その係数行を $(u,i)$ の順に並べた行列を $\widehat P$ とする。
$Q=PAP^t$ に対して
\[
 G(Q)_{(u,i),(v,j)}=
 \begin{cases}\left[x_i^2\right]Q_{uv},&i=j,\\
 \left[x_ix_j\right]Q_{uv}/2,&i\ne j\end{cases}
\]
を定めると、同符号の全incidenceから
$G(Q)=\widehat P A\widehat P^t$ となる。
混合直交性だけではこの別々の対称性は得られない。
この係数の偏極が、二つの復元法の共通の出発点である。

\subsection{四次の場合：支配性と残余符号}
正規化チャートは百五十三変数、六十九方程式、接次元八十四である。
同じ seed $17029$ の Hensel 分枝上で
\[
 dQ|_{T}=D_t-D_zJ_z^{-1}J_t
\]
の十八次小行列式は $101$ 進付値二、入力精度は $101^{32}$ である。
$J_z$ は単元で、逆行列の更新を最終補正座標で再確認しているため精度損失はない。
従って十八個の Gram 係数は独立に動く。
任意の非超楕円種数三 Prym 対は、基点を持たない $|K_C\eta|$ の二つの切断の積から
対称二次行列表示を持つので、これが $\mathcal R_3$ への支配性を証明する。

一般に $G(Q)$ は可逆で、関係枠を変えて
\[
 P(x)=\begin{pmatrix}x^t&0&0\\0&x^t&0\end{pmatrix},\qquad A=\operatorname{diag}(G,1)
\]
とできる。一次直交 syzygy 空間は $K=K_6\oplus K_3$、$\dim K_6=6$、$\dim K_3=3$ である。
$B$ の四行は $\operatorname{Gr}(4,K)$ の点 $W$ を与え、
三つの交代形式について同時等方的でなければならない。
標準化に残る $\tau(c,\ell)=(c,-\ell)$ を必ず商する。
この符号を忘れると、逆構成に現れる二価性を余分な被覆と取り違える。

pairing を正規化しないチャートの次元は八十五であり、
固定 Gram ファイバーの商は $85-18-49-16=2$ 次元である。
$A_{77}=1$ の八十四次元チャートでは、Gram を固定して動ける関係枠は四十八次元である。
八十四から四十九を差し引くことは、同じスカラーを二度消している。

\subsection{楕円線織ファイバーの逆構成}
$D_i=\operatorname{diag}([e_i]_\times,[e_i]_\times)$ とし、
\[
 \Omega_{ij}=D_i^tG^{-1}D_j-D_j^tG^{-1}D_i,\qquad
 E_Q=V\bigl(\operatorname{Pf}(z_1\Omega_{23}+z_2\Omega_{31}+z_3\Omega_{12})\bigr)
\]
とする。一般に $E_Q$ は滑らかで、その上の核束 $\mathcal K_Q$ は階数二である。
\[
 (X_Q\text{ の有限配置開部分})/\langle\tau\rangle
 \dashsim\mathbb P_{E_Q}(\mathcal K_Q).
\]
必要な条件 $\operatorname{rank}(W\to K_3)=3$ は同じ Hensel 点で証明されている。
有限配置条件は $W\cap K_3=0$ を強制する。非零な交わりがあれば、
$B$ の一行が $\ell(x)e_7$ となり、直線 $\ell=0$ 全体が階数低下集合に入るためである。

逆構成では $k=W\cap K_6$ を取り、
\[
 F_k=\{v\in K_6:k^t\Omega_{ij}v=0\text{ for all }i<j\}
\]
を作る。滑らかな三次曲線が直線を含まないことから、
$k\subset\ker\Omega(z)$ を満たす $[z]\in E_Q$ は一意で、$\dim F_k=4$ である。
$W$ は $F_k\to K_3$ のグラフとなり、三次元商 $F_k/k$ 上で
指定された可逆写像を $\bigwedge^2\phi$ に持ち上げる。
$\bigwedge^2:\operatorname{GL}_3\to\operatorname{GL}_3$ の核は $\{\pm1\}$ なので、
二つの持上げが $\tau$ でちょうど同一視される。

一般の平面四次曲線に自己同型がないことと、対称最小分解の一意性により、
粗 Prym ファイバーにさらに有限商は生じない。
ただし、上の階数二核束は一般関数体上で実際のベクトル束へは降りず、
その射影化は非零Brauer classを持つ二次曲線束へ降りる。
この算術的な違いを第\ref{subsec:notes-brauer-v27}節で記述する。

\subsection{六次の場合：有理式による再構成}
九行七列の $\widehat P$ は階数七、混合方程式 $PAb^t=0$ の十八行二十一列の行列は階数十八を持つ。
いずれも同じ Hensel 点で単元小行列式を持つので、非空開条件である。
$G(Q)$ の可逆な主七次部分行列 $G_I$ を選び、
\[
 A_0=G_I,\qquad \widehat P_0=G(Q)_{*,I}G_I^{-1}
\]
とすると、$G(Q)=\widehat P_0A_0\widehat P_0^t$ である。
他の階数七の分解との差は $\operatorname{GL}_7$ だけである。
次に三次元核 $K_{P,A}$ の基底を取れば、$B$ も行変換を除いて復元される。

この手順は像の判定法にもなる。任意の三行三列の対称二次式行列に対して、
まず $\operatorname{rank}G(Q)=7$、混合核の次元三を調べ、
その核上の三つの交代形式が消えることを検査する。
これに既存の開階数と有限配置条件を加えると全incidenceを再構成できる。
階数七だけから残りの九つの同符号方程式が自動的に成立すると考えてはならない。
どの既約成分が得られるかも、指定した $Z_1$ との照合が必要である。

平面 Prym データは対称最小分解から $Q$ を合同変換まで復元する。
さらに以下の一意性により、$H$ を指定せずに抽象対 $(C,\eta)$ から復元できる。
従って $\mathfrak Z_1^+\dashrightarrow\mathcal R_{10}$ は七次元の像への双有理写像であり、
全二十一点配置と特定の $D_{\rm mov}\in|K_C\eta|$ も戻る。

\paragraph{再利用できる平面直線束の一意性.}
標数零の代数閉体上の滑らかな平面曲線 $C$ の次数を $d\ge4$ とする。
次数 $d$ 以下で三つ以上の切断を持つ線束は $H=\mathcal O_C(1)$ だけであり、
Brill--Noether軌跡 $W_d^2(C)$ は被約な一点である。
これはNoether理論の特別な場合である\cite{HartshorneNoether,CoppensKatoPlaneSystems}が、
本研究で必要な証明は次のように短く閉じる。

固定因子を除いた線束 $M'$ の次数を $n\le d$ とし、一般の被約除数を $D$ とする。
随伴公式とRiemann--Rochにより、次数 $d-3$ の平面形式の評価階数は
\[
 \operatorname{rank}(\mathrm{ev}_D)=n+1-h^0(M')\le n-2.
\]
$n\le d-2$ なら、各点を、他の $n-1$ 点を一つずつ通る直線の積で分離でき、階数は $n$ となる。
$n=d-1$ で非共線なら、一点 $p$ を避けて残りの二点を結ぶ直線と、
残る $d-4$ 点を通り $p$ を避ける直線の積で全点を分離できる。
$d=4$ でも三つの非共線点を一次式で分離する証明そのものである。
共線なら階数は $d-2$ なので、いずれも上界に反する。
従って $n=d$ である。非共線な $D$ は非共線な $d-1$ 点を含むので再び矛盾し、
$D$ は直線上にある。B\'ezoutにより $D=C\cap\ell$、従って $M'=H$ となり、
次数から元の固定因子も零である。

$[H]$ における $W_d^2(C)$ の接空間は、Serre双対性によって乗法
\[
 H^0(H)\otimes H^0(K_C-H)\longrightarrow H^0(K_C),
 \qquad S_1\otimes S_{d-4}\longrightarrow S_{d-3}
\]
の像の零化空間である。この乗法は全射なので接空間は零となり、Nakayamaの補題が被約性を与える。
また、平面曲線の抽象同型は $H$ を保つため、射影変換へ一意に延長する。
六次で $H^3=K_C$ の根だけを数えると失われる情報は、切断数 $h^0(H)=3$ である。
この復元は曲線を保持しており、Prym多様体だけからのTorelliを意味しない。

\paragraph{固有ブロックの順序を戻す規則.}
署名 $(3,3)$ の有限配置開部分では、$P$ は平面全体で階数三、
$B$ はちょうど三固定点で階数低下する。
従って、二つのブロックのうち全点で最大階数を持つものを $P$ として選べる。
pairingのスカラー正規化を戻した枠付き開部分を $\mathcal U$、
順序を保つ枠・座標・pairing変更群を $\mathcal G_{\rm ord}$ とする。
$s(P,B,A)=(B,P,A)$ は $J$ を $-J$ に移し、全incidenceを保つ。
$\mathcal U\cap s\mathcal U=\varnothing$ であり、
\[
 (\mathcal U\sqcup s\mathcal U)/(\mathcal G_{\rm ord}\rtimes\langle s\rangle)
 \dashsim \mathcal U/\mathcal G_{\rm ord}
 \dashsim\mathfrak Z_1^+.
\]
交換は別の枠付き成分を動かし得るので、順序を忘れる商はこの合併の商として定義する。
一つの成分に交換が作用すると無条件に仮定しない。
ここで不要になったのは固有ブロックの順序であり、指定反対合そのものの一意性と
$Z_1$ の一般の全外部自己同型群は未判定である。

\subsection*{六次Prym像の再利用可能な標準形}

階数七の係数Gramの核は二次元pencilである。
行列式が相異なる三根を持つ開部分では、これを対角の跡零pencilに正規化できる。
正規化後の $P$ と三次元の一次直交syzygy基底 $R$ は
論文Vの「九本の余因子方程式」で与えた固定行列となり、
全ての変形情報は十九変数の対称形式 $A$ に移る。
負ブロックは $B=RA^{-1}$ であり、残る条件は
\[
 \begin{aligned}
 (A^{-1})_{14}&=(A^{-1})_{15}=(A^{-1})_{16}=0,\\
 (A^{-1})_{23}&=(A^{-1})_{24}=(A^{-1})_{26}=0,\\
 (A^{-1})_{35}&=(A^{-1})_{36}=(A^{-1})_{45}=0.
 \end{aligned}
\]
である。添字は $0,\ldots,6$。
従って、後続の計算を百五十三変数・七十二方程式のチャートから、
十九変数・九余因子方程式のチャートへ縮小できる。
この正規化に必要な三根の順序付けは、最後に $\mathfrak S_3$ で除く。

$A_0$ は $(0)$ の成分が $1$ で、$(1,3),(2,5),(4,6)$ を交換する行列である。
$\det A_0=-1$。
上記の順序の九余因子の $(p,m,i,r,o,l,b,g,f)$ に関するJacobi行列は $I_9$ である。
従って局所計算には $(s-1,a,c,d,e,h,j,k,n,q)$ の十変数を自由変数として用いられる。
これは九個の変数を形式的・エタール局所的に消去できるという主張であり、
大域的な有理パラメータ表示を与えるものではない。
また $A_0$ ではトーラスの安定化群が一般点より大きいため、
十次元チャートの滑らかさだけから粗商の基点の滑らかさを推論しない。

正則核pencilを持つ階数七Gram軌跡は、十四次元のpencil軌道上の
十九次元線形空間族の開部分であり、滑らかな三十三次元空間である。
九方程式は既存成分に沿う一般点で余次元九を持つ。
選択した十次元成分の有効安定化群は
$\mathbb G_m^3\rtimes\mathfrak S_3$ であり、その商は七次元である。
全余因子軌跡の既約性や他の成分は未判定。
成分の選択は既存の整数基点 $B_-$ と同定した $A_0$ により固定する。

\paragraph{標準的な直線とPrym-canonical切断.}
核pencilの三つの特異元から、両三次元空間の固有直線分解を取り戻せる。
対応する三つの階数一テンソルが張る三次元空間 $E$ は根の順序に依存せず、
$\ell=E/K\subset (U\otimes V)/K$ は標準的な直線となる。
これは標準形の $\langle e_0\rangle$ である。
$v\in\ell\setminus\{0\}$ に対する $t=PAv$ の余核像は
$H^0(C,K_C\otimes\eta)$ の標準的な一次元部分空間を定める。
十八点除数は直接
\[
 I_{D_{\rm mov}}=(\det Q,\operatorname{adj}(Q)t)
\]
で復元される。
その最小分解の次数は、生成元が $(5,5,5,6)$、syzygyが $(7,7,7)$。
三つの五次生成元は固定三角形の三頂点で可逆な対角評価行列を与える。
従って $h^0(I_D(5))=3$、$h^0(I_Z(5))=0$。
固定三点は核pencilだけから、移動十八点はこの標準切断から得られる。

補助六次曲線は、標準形で $v=e_0$ と正規化し、$s=v^tAv$ と置くと
\[
 \det A\det(RA^{-1}R^t)
 =s\det Q-t^t\operatorname{adj}(Q)t
\]
から得られる。この零曲線は $v$ の非零定数倍に依存せず、$Q$ と標準直線から計算できる。

\paragraph{次に調べる価値のある問題.}
十次元の余因子成分と三次元トーラスの表示により、有理性・不変量・分岐の計算が具体化した。
当面は九余因子方程式の大域的な消去や低次の不変量を調べることができるが、
形式局所解を有理解と取り違えないことが必要である。
全monodromyと反対合の標識忘却には、標準形が与える固定三角形以外の情報が要る。
また、九条件を単なる階数七条件から自動的に導いてはならない。
階数七・正則pencilの三十三次元軌跡と、本研究の二十四次元枠付き像は異なる。

\subsection{Hermite--Pfaffian恒等式と二次同種写像}
\label{subsec:notes-hermite-v27}

四次の対称表示を
\[
 Q=\begin{pmatrix}a&b\\b&c\end{pmatrix},\qquad
 a=x^tA_0x,\quad b=x^tA_1x,\quad c=x^tA_2x,
 \qquad G=\begin{pmatrix}A_0&A_1\\A_1&A_2\end{pmatrix}
\]
と書き、二次曲線のnetを $N=\langle a,b,c\rangle$ とする。
双対平面上の普遍階数二部分束 $\mathcal S$ に対する制限
\[
 r_N:N\otimes\mathcal O\longrightarrow\operatorname{Sym}^2\mathcal S^*
\]
の行列式を $H_N$ と書く。これは三次式であり、$H_N(z)=0$ は
netの非零な二次曲線が直線 $\ell_z$ を含むことと同値である。
具体的には $x=(z_3u,z_3v,-z_1u-z_2v)$ を代入し、
$a,b,c$ の係数を $u^2,uv,v^2$ の順に並べた行列 $R_N$ を使って
$H_N=\det(R_N)/z_3^3$ と計算できる。
分母は消え、全双対平面上の式となる。

既存の外積・Pfaffianの規約では、正確な恒等式は
\begin{equation}\label{eq:notes-hermite-pfaffian-v27}
 2\det(G)\operatorname{Pf}\bigl(z_1\Omega_{23}+z_2\Omega_{31}+z_3\Omega_{12}\bigr)
 =H_N(z)
\end{equation}
である。証明は一つの直線での相補小行列式に帰着する。
$M=G^{-1}$ とすると
\[
 \operatorname{Pf}\Omega_{12}
 =\det M_{(1,3,6),(5,3,6)}-\det M_{(2,3,6),(4,3,6)}.
\]
Jacobiの公式により、$\det(G)$ を掛けた右辺は
\[
 \det\begin{pmatrix}
 (A_0)_{11}&(A_1)_{11}&(A_2)_{11}\\
 (A_0)_{12}&(A_1)_{12}&(A_2)_{12}\\
 (A_0)_{22}&(A_1)_{22}&(A_2)_{22}
 \end{pmatrix}=\tfrac12 H_N(0,0,1)
\]
となる。中央の係数が二倍になることが定数 $2$ の由来である。
座標変更 $x=gx'$ では両辺とも
$F(z)\mapsto\det(g)^3F(g^{-t}z)$ と変わるので、全座標で恒等式が成立する。
従って補助曲線 $E=E_Q$ は、古典的なnetのHermite曲線そのものである
\cite[著者公開原稿 \S3.2.2, Exercise 3.12]{DolgachevCAG}。
その同型類は、$Q$ の三成分の基底よりもnet自体に依存する。

判別式三次曲線を
\[
 \Delta_N=V\bigl(\det(\lambda_0A_0+\lambda_1A_1+\lambda_2A_2)\bigr)
 \subset\mathbb P(N)
\]
とする。$E$ と $\Delta_N$ が滑らかなら、直線成分をそれを含む退化二次曲線へ送る写像
\[
 \pi_N:E\longrightarrow\Delta_N
\]
は連結有限エタール二重被覆である。
実際、判別式の滑らかさにより退化二次曲線の階数は常に二で、二直線は相異なる。
一つの直線が独立な二つのnetの元に含まれると、そのpencilは全て退化し、
$\Delta_N$ に直線成分を生じるため、像は一意に定まる。
局所的には階数二の非退化二次形式の等方直線を選ぶエタール二重被覆となる。

被覆のtrace零線束は対称表示
\[
 0\longrightarrow\mathcal O_{\mathbb P(N)}(-2)^3
 \xrightarrow{\lambda_0A_0+\lambda_1A_1+\lambda_2A_2}
 \mathcal O_{\mathbb P(N)}(-1)^3\longrightarrow i_*\theta_N\longrightarrow0
\]
の余核である。対称性から $\theta_N^2=\mathcal O_{\Delta_N}$、
コホモロジーから $h^0(\theta_N)=0$ となる。
同定を直接確認するには、$\mathscr K=\ker(\mathcal O^3\to\mathcal O(1)^3)$、
$\mathscr V=\mathcal O^3/\mathscr K$ と置く。
等方直線被覆の判別式線束は
$\det\mathscr V\otimes\mathcal O(-1)=\mathscr K^*\otimes\mathcal O(-1)$ であり、
上の余核 $\theta_N$ と一致する。従って
\[
 (\pi_N)_*\mathcal O_E=\mathcal O_{\Delta_N}\oplus\theta_N,
 \qquad
 \operatorname{Pic}^0(E)
 \simeq\operatorname{Pic}^0(\Delta_N)/\langle\theta_N\rangle.
\]

対称最小分解の一意性により、$(\Delta_N,\theta_N)$ は一般Prym対の不変量である。
これが支配的有理写像
\[
 \Psi:\mathcal R_3\dashrightarrow Y_0(2),\qquad
 (C,\eta)\longmapsto
 (\operatorname{Pic}^0(\Delta_N),\langle\theta_N\rangle)
\]
を与える。支配性は後述の試験族の非定数な $j$ で確認できる。
Jacobianを用いるので原点の選択を要しない。
この二次同種写像から、種数三Prym対の二次元Prym多様体に楕円因子があるとは結論しない。
また、核束までnetだけに依存するという主張は含めない。

\subsection{net写像の相対有理性と降下で残るもの}
\label{subsec:notes-relative-rational-v28}

$\Psi:\mathcal R_3\dashrightarrow Y_0(2)$ は
Dolgachev\cite[Remark~4.1]{DolgachevRationality2008}の古典的なnet写像である。
本研究で得た相対的な結論は
\[
 \mathbb C(\mathcal R_3)=\mathbb C(Y_0(2))(u_1,\ldots,u_5)
\]
という、指定された底を保つ純超越表示である。公表稿VIに全証明を置く。
再利用すべき点は、一般netの非自明な安定化群を消さずに扱う方法である。

一般のnetの射影安定化群は
$G=(\mathbb Z/3)^2\rtimes C_2$ である。
$\operatorname{SL}_3$ 内の逆像 $\widetilde G$ を、その一意のSylow三部分群による
符号指標 $\chi$ でねじると、線形持上げは
三乗根の対角群と $S_3$ の半直積 $G(3,3,3)$ となる。
その不変式環は
\[
 k[p,q,r],\qquad \deg(p,q,r)=(3,3,6)
\]
であり、分裂形の生成元は
$p=x^3+y^3+z^3$、$q=xyz$、$r=x^3y^3+y^3z^3+z^3x^3$ である。
非定数群の形に対しても、次数三の二次元空間と次数六の一次元補空間を
基礎体上で選ぶことにより同じ多項式環が降りる。
従って $\mathbb P^2/G\simeq\mathbb P(1,1,2)$ は分裂している。

さらに $V'=V\otimes\chi$ とすると
\[
 \operatorname{Sym}^2V^*\simeq V'\otimes M,\qquad
 M=\operatorname{Hom}_{\widetilde G}(V',\operatorname{Sym}^2V^*),\quad\dim M=2.
\]
二つの部分表現を個別に降ろす代わりに、このHom空間を使う。
$\mathbb P(V'\oplus V')\dashrightarrow\mathbb P(V')$ の開部分は
階数三のHom束であり、中心 $\mu_3$ はこの束上で相殺される。
自由作用開部分で商を取れば真のベクトル束となるため、二次曲線空間の商は五次元有理である。
この議論は定数群の標準形を基礎体へ無条件に移す議論ではない。

モジュラー関数体上のnet代表は
\[
 \langle y^2+2xz,\ 2yz,\ -x^2+2xz-2y^2+uz^2\rangle
\]
と取れる。判別式のWeierstrass表示と $j$ は
\[
 Y^2=X(X^2+6X+9+u),\qquad
 j=\frac{1728(u-3)^3}{u(u+9)^2}.
\]
唯一の非零有理二捩れ $(0,0)$ と次数三の $j$ 写像により、$u$ は $Y_0(2)$ の座標となる。
入力の二次曲線は $\mathbb P(N)$ 内にあり、その双対を評価写像で引き戻すと四次Prym曲線になる。

底の関数体の有理性は、一般Hermite曲線やPrym-canonical二次曲線の分裂とは独立である。
有限安定化群による降下が残るため、周期三・指数二・積の指数六と両立する。
また $\Psi$ 全体が正則な射影空間束であるという結論ではない。


\subsection{一般Prymファイバーの非零Brauer障害}
\label{subsec:notes-brauer-v27}

$K=\mathbb C(\mathcal R_3)$ とし、曲線の自己同型が自明な開部分上の一般対を
\[
 C/K,\qquad 0\ne\eta\in\operatorname{Pic}^0_{C/K}(K)[2]
\]
とする。$\eta$ はまずPicardスキームの点であり、実際の $K$ 上の線束とは限らない。
その降下障害を
\[
 \alpha_{\mathrm{Prym}}=\delta_C(\eta)\in\operatorname{Br}(K)[2]
\]
と書く。これは元の十一点族のpotential障害 $\beta$ とは別の対象である。
Brauer classと射影線形群の降下cocycleの対応は
\cite[Tags 03R1, 0A2J]{StacksBrauerV27}を用いる。

$V_{\rm can}=H^0(C,\omega_C)$ は真の三次元 $K$ ベクトル空間だが、
$W=H^0(C_{\overline K},\omega_C\eta)$ は $\alpha_{\rm Prym}$ による階数二のtwisted空間である。
その射影化はBrauer--Severi二次曲線 $D/K$ へ降りる。
$\eta^2\simeq\mathcal O_C$ を選ぶ $\mu_2$-gerbe上では、
非自明慣性は $W$ に $-1$ として作用する。
乗法 $\operatorname{Sym}^2W\to\operatorname{Sym}^2V_{\rm can}$ は $Q$ を与え、
係数Gram形式を $W\otimes V_{\rm can}^*$ 上の形式と解釈できる。
この規約で余核は正確には
$(\omega_C\eta)\otimes\det(W)^{-1}$ であるが、$\det(W)$ は慣性が自明な $K$ 上の線であり、
その自明化による調整は障害を変えない。

一次直交syzygyの空間は、Gram形式で値を双対化すると
\[
 T\simeq W^*\otimes\bigwedge^2V_{\rm can}
\]
となるので、慣性は $T$ にも $-1$ として作用する。
三つの交代形式は自然な写像
$\bigwedge^2T\to\bigwedge^2V_{\rm can}$ を与える。
そのパラメータ平面は
$\mathbb P((\bigwedge^2V_{\rm can})^*)\simeq\mathbb P(V_{\rm can})$ であり、$K$ 上分裂している。
従って $E$ は $K$ 上の平面三次曲線へ降りる。
その上の階数二Pfaffian核にも慣性は $-1$ として作用するので、
射影化はBrauer class $\alpha_{\rm Prym}|_E$ を持つ二次曲線束 $S\to E$ へ降りる。
残余符号商を取ったファイバーの逆構成はこの降下と可換である。
同じ次数二のBrauer classを持つ二次曲線は同型なので、一般点上で
\begin{equation}\label{eq:notes-prym-product-v27}
 S\dashsim_K E\times_KD
\end{equation}
を得る。これは一般ファイバーの双有理モデルであり、束全体の積表示ではない。

\paragraph{非消滅を検出する支配族.}
$F=\mathbb C(s,t)$ 上の二次曲線
$D_0:Z^2-sX^2-tY^2=0$ は非分裂で、独立変数を添加しても非分裂である。
$s$ 進付値を使うと、剰余体上の $Z_0^2=tY_0^2$ から $Z_0=Y_0=0$ となり、
残る $sX^2$ の付値一を打ち消せないことが直接の証明になる。
十八個の独立係数からなる一般の三二次式で $f:\mathbb P^2\to\mathbb P^2$ を作り、
$C=f^{-1}(D_0)$ とする。一般には基点がなく、$C$ は滑らかな四次曲線である。
$D_0$ の幾何学的次数一のPicard classを $\ell$ とし、
$M=f_C^*\ell$、$\eta=M\omega_C^{-1}$ と置く。
$\ell^2=\mathcal O_{D_0}(1)$ から $\eta^2=\mathcal O_C$、
対称表示の分解から $h^0(M)=2\ne3=h^0(\omega_C)$ なので $\eta\ne\mathcal O_C$ である。
自然性により
\[
 \delta_C(\eta)=\delta_{D_0}(\ell)=[D_0]=(s,t)\ne0.
\]
代数閉包上ではこの族は全ての一般二行二列対称二次式表示を動かし、
$\mathcal R_3$ を支配する。従って一般の $\alpha_{\rm Prym}$ は非零で、位数・指数ともに二である。
ここで非分裂性を保つ純超越拡大は $F\subset F(\text{十八係数})$ であり、
$K$ からこの体への拡大が純超越であるとは仮定していない。

\paragraph{曲線へ制限しても分裂しない理由.}
$E$ は分裂平面内の三次曲線なので、次数三の零cycleを持つ。
もし $\alpha_{\rm Prym}|_{K(E)}=0$ なら有理写像 $E\dashrightarrow D$ が存在し、
滑らかな曲線から固有曲線への写像なので全体に延びる。
次数三の零cycleを押し出すと $D$ 上に次数三の因子ができ、
その線束と $\omega_D$ の積は次数一になる。
種数零のRiemann--Rochから $D(K)\ne\varnothing$ となり矛盾する。
従って一般の核束の射影化は、$E$ 上の真の階数二ベクトル束の射影化にはならない。
一方、$S$ の一般点はその二次曲線上の点を与えるので、$\alpha_{\rm Prym}|_{K(S)}=0$ となる。

\subsection{補助曲線の指数三と最小多重切断次数六}
\label{subsec:notes-index-v27}

曲線や固有多様体の指数は、全閉点の剰余体次数の最大公約数とする。
種数一曲線の周期は、そのJacobianに対するtorsor class
$[E]\in H^1(K,\operatorname{Pic}^0(E))$ の位数とする。
一般の $E/K$ は原点を持たず、周期と指数はともに三である。
上界 $\operatorname{ind}(E)\mid3$ は平面直線切断から従う。
下界には、単なる一つの特殊化でなく、次の\emph{支配的な}付値族を使う。
超越元 $r$ と離散付値パラメータ $s$ に対して
\[
 M_s(u,v,w)=\begin{pmatrix}
 ru&sw&sv\\sw&rsv&su\\sv&su&rs^2w
 \end{pmatrix}
\]
という二次曲線netを取ると、判別式とHermite曲線の式は
\begin{align*}
 \det M_s&=-rs^2\left(u^3+sv^3+s^2w^3-\frac{r^3+2}{r}suvw\right),\\
 E_s&:\quad r^2(s^2z_1^3+sz_2^3+z_3^3)+(r^3-4)sz_1z_2z_3=0.
\end{align*}
$t^3=s$ を入れると、$M_s$ は座標変更 $(u,v,w)=(U,V/t,W/t^2)$ と
$\operatorname{diag}(1,t,t^2)$ による合同変換で $M_1$ から得られる。
二つの三次曲線の滑らかさを保証する因子はそれぞれ
$(r^3-1)^2(r^3+8)$ と $(r^3-1)(r^3+8)^2$ であり、$r$ が超越元なら非零である。

支配性の再現可能な判定は次である。$M_1$ の三係数二次式を $(q_0,q_1,q_2)$ とし、
平面座標 $g\in\operatorname{GL}_3$ とnet基底 $h\in\operatorname{GL}_3$ を加えた
\[
 (r,g,h)\longmapsto h\,(q_0(gx),q_1(gx),q_2(gx))^t
\]
を考える。十九変数から十八係数への接行列は $r=2,g=h=I$ で階数十八を持ち、
一つの十八次小行列式は $-462422016$ である。
これにより全順序付きnetの空間、従って一般対称表示を通じて $\mathcal R_3$ を支配する。
$g,h$ の独立な成分を $\mathbf a$ とすれば
$K\hookrightarrow F=\mathbb C(r,\mathbf a)((s))$ であり、$E_F\simeq E_s$ となる。
この支配性は、指数の計算をモジュライの一般点へ戻すために不可欠である。

有限拡大 $L/F$ の次数が三で割れなければ、分岐指数 $e=v_L(s)$ も三で割れない。
三座標の付値を $a,b,c$ とすると、三つの立方項の付値
\[
 2e+3a,\qquad e+3b,\qquad3c
\]
は三を法として相異なる。混合項の付値 $e+a+b+c$ はこの三つの平均なので、
唯一の最小付値より真に大きく、方程式は成立しない。
零の座標があっても同じ付値比較で矛盾する。
従って全閉点の次数は三の倍数である。
一方、$t^3=s$ 上には $[t^{-2}:-t^{-1}:0]$ があるので $\operatorname{ind}(E_F)=3$ である。
基礎体拡大で指数は元の指数を割るため、$\operatorname{ind}(E)=3$ を得る。
種数一曲線の周期は指数を割り、周期一ならJacobianの自明torsorとして有理点を持つので、周期も三である。

式\eqref{eq:notes-prym-product-v27}の両射影から
$\operatorname{lcm}(\operatorname{ind}E,\operatorname{ind}D)=6$ は積の指数を割り、
零cycleの外積から積の指数は $3\cdot2=6$ を割る。従って
\[
 \operatorname{ind}(S)=6.
\]
滑らかな射影曲面の点吹上げは指数を変えないので、同じ値は全ての滑らかな射影双有理モデルに成立する。

さらに、指数六の定義だけに依存せず、任意の非空開部分内に次数六の閉点を作れる。
$U\subset E\times D$ の像の補集合を避ける一般 $K$ 直線で $E$ を切ると、
指数三により一つの次数三の閉点 $x$ を得る。
$L=\kappa(x)$ 上でも $D_L$ は非分裂である。分裂すれば $D$ に次数三の零cycleができるからである。
$D_L$ を $U_x$ の補集合を避ける一般 $L$ 直線で切ると、一つの次数二の閉点 $y$ を得る。
その合成次数は $[\kappa(y):K]=2\cdot3=6$ である。
双有理写像の定義される開部分を通じて、同じ構成を一般ファイバーへ戻せる。
正則な一般ファイバーの閉点は固有な滑らかモデルへのLang--Nishimuraにより指数六に従うので、
全ての多重切断次数は六の倍数であり、最小次数は六である。
この結論は正則な商開部分または滑らかなモデルに関する。
任意に収縮した特異境界に入る切断まで、同じ双有理議論で禁止してはならない。

\subsection{Pl\"ucker二次曲線と内在的二次拡大}
\label{subsec:notes-quadratic-splitting-v28}

一般関数体 $K=\mathbb C(\mathcal R_3)$ 上のweight一のsyzygy空間 $T$ と
$\Omega:\bigwedge^2T\to B$、$\dim T=6$、$\dim B=3$ から
\[
 c_\Omega:\bigwedge^3T\to T\otimes B,\qquad
 uvw\mapsto w\otimes\Omega(u,v)-v\otimes\Omega(u,w)+u\otimes\Omega(v,w)
\]
を作る。一般に階数十八、核 $\Pi$ はweight一の階数二空間である。
従って $D_\Pi=\mathbb P(\Pi)$ はPrym Brauer class
$\alpha_{\rm Prym}$ を持つ二次曲線であり、$D_\alpha$ と $K$ 同型となる。
同型の選択まで標準的であるとは主張しない。

Pl\"ucker関係式の制限は一つの二次式を生成し、
\[
 Y=\operatorname{Gr}(3,T)\cap D_\Pi=\operatorname{Spec}L
\]
は次数二のエタール因子となる。
$D_\Pi$ は非分裂なので $L/K$ は真の二次体であり、$\alpha_{{\rm Prym},L}=0$ である。
これは二つの共通等方三平面を区別する内在的な分裂体である。
一方、$\operatorname{per}(E_L)=\operatorname{ind}(E_L)=3$ は残り、
$S_L\dashsim E_L\times\mathbb P^1$、$\operatorname{ind}(S_L)=3$ となる。

基礎体上で二次代数を計算するには、次数六の中心単純代数 $\operatorname{End}(T)$ の中で
\[
 \mathfrak a_\Omega=\{J:\Omega(Ju,v)+\Omega(u,Jv)=0\}
\]
という一次元空間を解けばよい。
その非零元は $J^2=\delta I$、$\operatorname{Trd}(J)=0$ を満たし、
$L=K[J]=K[t]/(t^2-\delta)$ である。
$rI+sJ$ は $\Omega$ を $r^2-\delta s^2$ 倍するため、
$\operatorname{Res}_{L/K}\mathbb G_m/\mathbb G_m$ は $S\to E$ に忠実に作用する。
その固定部分は $E_L\to E$ という次数二の有限エタール多重切断である。
二次拡大上で核束は次数 $-3$ の二つの線束に分かれる。

試験族 $G(t)$ では
\[
 q_t(a,b)=(t-4)a^2+(t-4)ab-b^2,\qquad \delta=t(t-4).
\]
従って $\sqrt{\det G}=\sqrt{t-4}$ では正しい二次代数を得ない。
$D_\Pi$ を射影直線として座標化できるのはBrauer classの分裂体上である。
その体 $k$ でも二次因子が分裂し得るので、一般式は
$k[t]/(t^2-\operatorname{disc}q)$ と書き、常に体だとは仮定しない。
計算入口は \nolinkurl{isotropic_plucker_v28.py} である。


\paragraph{次の判定に残るもの.}
一般Brauer classの非消滅、$E$ の原点不存在、指数三・六は確定している。
内在的な二次分裂体も確定した。残るのは境界での分岐と相対コンパクト化の降下、全八次元空間の有理性、
および指定反対合自体を忘れたモジュライの幾何である。
元の十一点族の相対potential障害 $\beta$ の消滅は、この議論からは判定されない。


\subsection{補助曲線の節点と使い回せる長さ判定}
固定点 $x_0$ で $M=L(x_0)$、左核 $u=(0,b)$、$w_i=uL_i$ と置く。
全incidenceは $w_iAM^t=0$ を与え、被約な階数低下点の法線微分は
\[
 T_{x_0}\mathbb P^2\xrightarrow{\sim}(\ker M)^*
\]
である。従って $(w_iAw_j^t)$ は非退化であり、補助 Gram 行列式の二次初項を与える。
固定点は常に通常二重点である。

補助曲線の三偏微分に対する次数 $D=3d-5$ の Macaulay 行列で、
$p=3$ では $102$ 次小行列式の付値 $48$、$p=2$ では $205$ 次小行列式の付値 $63$ を得る。
入力・保証精度はともに $101^{256}$ である。
これにより指定次数の Jacobian 商の次元はそれぞれ三、五以下となる。
再利用できる一般補題は
\[
 \dim_k(S/I)_D\le f\le D
 \quad\Longrightarrow\quad
 \operatorname{length}\operatorname{Proj}(S/I)\le f
\]
である。飽和化して一般一次非零因子で割れば、標準次数付き環の Hilbert 関数の差が、
零になるまで正であることから従う。
一点の次数の商次元を、根拠なく射影スキームの長さと同一視しない。

既知の節点がこの長さを尽くすので、他の特異点はない。
節点数が $d-1$ より少ないため、B\'ezout によって幾何学的既約性を得る。
六次では三節点・正規化種数七、八次では五節点・正規化種数十六である。
非固定点はすべて滑らかな通常二重接触点になる。
補助 Gram 行列は節点を含めて corank 一なので、節点曲線上にも
非自明二捩れ線束 $\eta_B$ があり、連結な有限エタール二重被覆を定める。
正規化への引き戻しの非自明性は未決定である。

\subsection{固定三点のmonodromyと残る独立性}
$Z_1$ では固定三点への幾何学的 monodromy の射影が $S_3$ である。
基点の退化イデアルの飽和は $(xy,xz,yz)$ であり、
連結な枠変更群に属する境界安定化元が三巡回置換と互換を与える。
三点の有限エタール被覆が境界まで延びるため、全二十一点が被約な稠密開部分にもこの作用が残る。
ここから $S_3\times\{1\}$ の包含は導けず、非固定側との独立性は別の課題である。

\subsection{補助楕円曲線の変動とGrassmann境界}
公表稿Vの一変数行列 $G(t)$ は
\[
 \det G(t)=t-4,\qquad
 E_t:\ 2tz_1^3+z_2^3+tz_3^3+(t-8)z_1z_2z_3=0
\]
を与える。これは三次曲線の滑らかさの明示証明であるとともに、
次の補助不変量の変動を追うための小さな試験族になる。
\paragraph{試験族の二つの $j$ 不変量.}
Hermite--Pfaffian恒等式により $E_Q$ の基底不変性はnetの幾何から分かる。
上の試験族に対応する判別式曲線は
\[
 \Delta_t:\quad\lambda_0^3+4\lambda_1^3+t\lambda_2^3
 -(t+4)\lambda_0\lambda_1\lambda_2=0.
\]
$t=1$ では $E_t,\Delta_t$ ともに滑らかで、$\det G=-3\ne0$ である。
通常の $j(i)=1728$ の正規化では
\begin{align*}
 j(E_t)&=-\frac{(t-8)^3\bigl((t-8)^3-432t^2\bigr)^3}
 {2t^2(t-2)^3(t+16)^6},\\
 j(\Delta_t)&=\frac{(t+4)^3\bigl((t+4)^3+864t\bigr)^3}
 {4t(t-2)^6(t+16)^3}.
\end{align*}
Hesse形 $X^3+Y^3+Z^3+3hXYZ=0$ の公式
$27h^3(8-h^3)^3/(1+h^3)^3$ から得られる
\cite[\S2]{ArtebaniDolgachevHesse}。
いずれも非定数であり、$\mathcal R_3\dashrightarrow Y_0(2)$ の支配性を確認できる。
十八次元の行列空間上の有理関数の非定数性なので、元の四次Gram曲線が滑らかな
稠密開部分へ制限しても変わらない。
二つの $j$ が満たす次数二のモジュラー方程式は恒等式として再検査できるが、
被覆の存在とtrace零線束の同定は第\ref{subsec:notes-hermite-v27}節の幾何学的証明に依拠する。

\paragraph{三成分閉包と局所モデル.}
一般の $Q$ について、$X_Q$ は滑らかな楕円線織面 $R$ と二つの平面
$\Pi_\pm\simeq\mathbb P^2$ からなる単純正規交差曲面である。
二つの共通等方三平面を $U_\pm$ とすると
$\Pi_\pm=\{U_\pm\oplus\ell:\ell\in\mathbb P(K_3)\}$ である。
射影 $W\to K_3$ の階数二は、Pfaffian三次曲線に直線を強制するため存在しない。
階数一の点はちょうどこの二平面、階数三の点は滑らかな一つの既約開部分となる。

$K_6=U_+\oplus U_-$ に対する交差ブロックを $M(z)$ とすれば、
境界近傍で全方程式は正確に
\[
 M(\ell)d=0
\]
へ消去できる。
$\det M=0$ は滑らかな三次曲線なので、局所式は $fv=0$ となる。
これは成分や交差を点集合だけでなくスキームとして決める小さなモデルである。
$R$ は $E$ 上の分解可能な射影直線束で、二重曲線 $D_\pm$ は互いに交わらない二つの切断となる。
境界平面の $\ell$ と線織面の底の $z$ を同じ座標と仮定しない。

\[
 \deg N_{D_\pm/R}=0,\qquad
 \deg N_{D_\pm/\Pi_\pm}=9,\qquad
 \deg\mathcal T_X^1|_{D_\pm}=9.
\]
一般一次摂動は種数二十K3曲面を与え、その全空間には二つの二重曲線上に九個ずつ普通二重点が生じる。
この九点は一次平滑化切断の単純零点である。
二つの平面成分のイデアルを吹き上げる局所式は $ab=tc$ の $(a,t)$ のblow-upであり、
射影的な小さなcrepant解消を与える。
中心ファイバーは、中央の $R$ を保ち、両側の平面をそれぞれ九点で吹き上げた三成分のType II Kulikovモデルとなる。
法束は平滑化零因子を引くことで相殺する。
九点を吹き上げた各平面を、追加条件なしに有理楕円曲面とは呼ばない。

\paragraph{スクロールの分極と再利用可能な数値.}
$\det\mathcal K=\mathcal O_E(-2)$ なので、
直線部分空間をパラメータ化する $\mathbb P_E(\mathcal K)$ の超平面類の平方は六である。
GrassmannのPl\"ucker類については $\deg X_Q=38$、両平面の次数は一ずつ、
主成分 $R$ の次数は三十六である。
これら二つの分極を同一視しない。
一般関数体上で二平面を区別するには、前述の二次拡大 $L/K$ が必要であり、
Kulikov修正の粗モジュライ上への降下とBrauer classの境界での分岐は残る問題である。

